\subsection{Пружинки}
%%%%%%%%%%%%%%%%%%%%%%%%%%%%%%%%%%%%%%%%%%%%%%%%%%%%
%%%%%%%%%%%%%%%%%%%%%%%%%%%%%%%%%%%%%%%%%%%%%%%%%%%%
\z{Нить, перекинутая через блок, изначально не натянута. К свободному концу нити прикладывают силу $F$. На сколько сместится конец нити под действием этой силы? Определите эквивалентный коэффициент жесткости системы. Жесткости пружин известны и указаны на рисунке.
		\begin{center}%
 			\includegraphics[width=0.9\textwidth]{3.4.1}%
		\end{center}%
}
%
{\puan $\frac{k}{4}$;~\puan $\frac{k}{4}$;~\puan $\frac{4k}{17}$;~\puan $\frac{k}{5}$;~\puan $\frac{k}{6}$;~\puan $\frac{k}{19}$;}

%%%%%%%%%%%%%%%%%%%%%%%%%%%%%%%%%%%%%%%%%%%%%%%%%%%%
\z[{\begin{wrapfigure}[2]{r}{0.3\textwidth}%
		\vspace{-15pt}%
		\begin{flushright}%
 			\includegraphics[width=0.3\textwidth]{3.4.2}%
		\end{flushright}%
	\end{wrapfigure}%
	}]%
{Брусок массой $1.6 \, \text{кг}$ тянут по гладкому столу с помощью пружины жёсткостью $40 \, \text{Н/м}$, сообщая бруску ускорение $0.5 \, \text{м/с}^2$. Каково удлинение пружины?}
%
{$x=ma/k=2 \,\text{см}$}

%%%%%%%%%%%%%%%%%%%%%%%%%%%%%%%%%%%%%%%%%%%%%%%%%%%%
\z[{\begin{wrapfigure}[3]{r}{0.3\textwidth}%
		\vspace{-15pt}%
		\begin{flushright}%
 			\includegraphics[width=0.3\textwidth]{3.4.3}%
		\end{flushright}%
	\end{wrapfigure}%
	}]%
{В недеформированном состоянии две пружины имеют длину $l$ и $2l$ и жёсткость $k$ и $2k$ соответственно. На каком расстоянии от левой стенки будет покоиться небольшое тело, закрепленное между пружинами, если расстояние между точками крепления пружин к стенкам $4l$?}
%
{$L=5l/3$}

%%%%%%%%%%%%%%%%%%%%%%%%%%%%%%%%%%%%%%%%%%%%%%%%%%%%
\z{Резиновый шнур жесткостью $k=100\, \text{Н/м}$ сложили вчетверо. Какова жесткость получившейся системы?}
%
{$k'=16k= 1600\, \text{Н/м}$}

%%%%%%%%%%%%%%%%%%%%%%%%%%%%%%%%%%%%%%%%%%%%%%%%%%%%
\z[{\begin{wrapfigure}[3]{r}{0.15\textwidth}%
		\vspace{-30pt}%
		\begin{flushright}%
 			\includegraphics[width=0.15\textwidth]{3.4.5}%
		\end{flushright}%
	\end{wrapfigure}%
	}]%
{Два груза массой $3m$ и $m$ подвешены к потолку с помощью невесомой пружины и лёгкой нити. Система покоится. Определить ускорение (направление и модуль) каждого груза сразу после пережигания нити.}
%
{$a_\text{верх}=\frac{4g}{3}$ вниз; ~ $a_\text{ниж}=0$}

%%%%%%%%%%%%%%%%%%%%%%%%%%%%%%%%%%%%%%%%%%%%%%%%%%%%
\z[{\begin{wrapfigure}[3]{r}{0.15\textwidth}%
		\vspace{-5pt}%
		\begin{flushright}%
 			\includegraphics[width=0.15\textwidth]{3.4.6}%
		\end{flushright}%
	\end{wrapfigure}%
	}]%
{Какую силу нужно приложить к подвижному блоку, подвешенному на двух пружинах, имеющих коэффициенты жесткости $k$ и $2k$, чтобы сместить его вниз на расстояние $x$?}
%
{$F=\frac{8kx}{3}$}
%%%%%%%%%%%%%%%%%%%%%%%%%%%%%%%%%%%%%%%%%%%%%%%%%%%%
\z[{\begin{wrapfigure}[9]{r}{0.3\textwidth}%
		\vspace{+30pt}%
		\begin{flushright}%
 			\includegraphics[width=0.3\textwidth]{3.4.7}%
		\end{flushright}%
	\end{wrapfigure}%
	}]%
{Найдите общий коэффициент жёсткости системы пружин, изображённой на рисунке, если внешняя сила прикладывается к верхней платформе в вертикальном направлении. Лестница, на которую опираются пружины, бесконечна. Все платформы при сжатии пружин сохраняют горизонтальное положение и не касаются ступенек лестницы. Каждая из платформ, кроме самой верхней, опирается на две пружины. Коэффициенты жёсткости всех пружин одинаковы и равны $k$, оси всех пружин вертикальны.}
%
{$k\sqrt{5}$}